\section{Kaggle implementation: an interpretable decision tree}
\label{sec:kaggle-titanic-tree}

\begin{kaggleproject}{Titanic --- constrain tree depth and inspect learned rules}
\kagglesource{https://www.kaggle.com/datasets/yasserh/titanic-dataset}{Titanic Dataset}
\textbf{Question.} What rules does a shallow decision tree learn, and how does constraining depth change the bias--variance tradeoff?
\end{kaggleproject}

The code deliberately limits \texttt{max\_depth} and \texttt{min\_samples\_leaf}, then prints the top of the fitted tree as text. This makes the model auditable enough for a teaching example while still using the full preprocessing pipeline.

\textbf{Before running.} Predict the qualitative pattern you expect from the experiment before you look at the output or plot.

\begin{labmode}
Stage 3B: choose a constrained tree configuration and define what evidence would count as overfitting before opening the reference implementation.
\end{labmode}

\textbf{Part A --- fit a constrained tree and measure the train--validation gap.}

\lstinputlisting[style=mlpython,firstline=1,lastline=41]{all_experiments/lab11-titanic_tree.py}

\textbf{Part B --- inspect and visualize only the readable top of the fitted tree.}

\lstinputlisting[style=mlpython,firstline=42,lastline=62]{all_experiments/lab11-titanic_tree.py}


\begin{debugtask}
Set \texttt{max\_depth=None} and \texttt{min\_samples\_leaf=1}. Compare training and validation accuracy with the constrained tree. Identify the symptom of overfitting from the train--validation gap rather than repeatedly probing a final test set.
\end{debugtask}

\begin{labdeliverable}
Submit the constrained tree's top learned rules, a table comparing constrained and unconstrained tree training/validation performance, and a short diagnosis of overfitting based on the train--validation gap rather than test score alone.
\end{labdeliverable}
